John Quincy Adams proposed his apportionment method in 1832 during
the congressional debate that followed the 1830 Census.
Adams, a former president serving in the House, objected to Jefferson's
method on grounds that it systematically underrepresented small states.
His proposed remedy was the opposite extreme: always round up.
Under Adams, every state receives at least the ceiling of its
exact proportional quota, ensuring that no state loses representation
due to fractional rounding.

Adams's method was never adopted for federal congressional apportionment.
In the 1832 debate, Congress ultimately chose neither Adams nor Jefferson
but moved toward Webster's arithmetic mean rounding \citep{balinski1982fair}.
The political economy of apportionment reform consistently favoured
compromises between small-state and large-state interests,
and Adams's ceiling rounding --- which gave small states an extra seat
even when their exact quota was barely above an integer ---
was seen as excessive in the other direction from Jefferson.

Despite its lack of federal adoption, Adams is theoretically important
for three reasons.
First, it defines one extreme of the divisor method spectrum:
Adams is the unique divisor method that is most small-state-favoring
(just as Jefferson is the unique divisor method that is most
large-state-favoring).
The spectrum Adams $\succ$ HH $\succ$ Webster $\succ$ Jefferson
provides the theoretical context for understanding what Huntington-Hill's
choice of geometric mean rounding achieves --- it is a deliberate
position between the two extremes.
Second, Adams's maximisation of minimum per-capita representation
(proved in Section~\ref{sec:bias}) is a meaningful fairness criterion:
it ensures that the most underrepresented state is as well-represented
as possible.
Third, Adams is paradox-immune as a divisor method, providing a concrete
example that paradox immunity does not require near-neutral rounding ---
even ceiling rounding, which is aggressively biased toward small states,
preserves the monotone allocation property.

\subsection{Scope}

This paper proves Adams's small-state optimality (Section~\ref{sec:bias}),
establishes paradox immunity (Section~\ref{sec:paradox}),
computes the 2020 Census Adams apportionment and compares it to
Huntington-Hill and Webster (Section~\ref{sec:empirical}),
and documents the \textsc{bisect-apportion} implementation
(Section~\ref{sec:implementation}).
